$\Cl^{n}$ is one-dimensional (\claimref{std-algebra.grade-decomposition}), spanned by $e_1 \wedge \dots \wedge e_n$ for a basis $e_i$, and $\underline{f}$ preserves grade, so $\underline{f}(I) = \lambda I$ for a scalar $\lambda$ independent of $I$. Write $f(e_j) = \sum_i f_{ij} e_i$. By multilinearity and alternation of $\wedge$ (\claimref{std-algebra.exterior-isomorphism}), $f(e_1) \wedge \dots \wedge f(e_n) = \sum_{\sigma} f_{\sigma(1)1} \cdots f_{\sigma(n)n}\, e_{\sigma(1)} \wedge \dots \wedge e_{\sigma(n)} = \bigl(\sum_\sigma \operatorname{sgn}\sigma \prod_j f_{\sigma(j)j}\bigr) e_1 \wedge \dots \wedge e_n$, the Leibniz formula; so $\lambda = \det f$. Since $\underline{fg}$ and $\underline{f}\,\underline{g}$ both respect $\wedge$ and agree on $V$, they are equal, and applying them to $I$ gives $\det(fg) = \det f \det g$.

If $f$ is an isometry, $\Cl(f)$ is an algebra homomorphism (\claimref{std-algebra.functoriality}), so it maps the antisymmetrised product $[v_1, \dots, v_k] = v_1 \wedge \dots \wedge v_k$ to $[f v_1, \dots, f v_k]$; hence $\Cl(f)$ agrees with $\underline{f}$ on every blade, so on everything. For $V$ nondegenerate, $f$ preserves the Gram determinant of a basis, which is nonzero, so $(\det f)^{2} = 1$.
